\section{Splitting method overview}

The splitting method is based on a estimation method to calculate the
free energy difference between two highly overlapping canonical ensembles\cite{bennett1976efficient}.
Consider a system with a configuration space defined by coordinates $\{q_1, \ldots, q_N\}$,
on which we define a temperature-scaled energy function $U_0(q_1, \ldots, q_N) = \Phi_0(q_1, \ldots, q_N)/k_BT$.
With this energy, we can build the partition function for the ensemble:
$$ Z_0 = \int \exp[-U_0(q_1, \ldots, q_N)] dq_1\ldots dq_N $$
An exact calculation of $Z_0$ is often intractable, but given the same configuration space, the ratio
between two partition functions $Z_0/Z_1$ defined by energy functions $U_0$ and $U_1$ is tractable by Monte Carlo methods.

To calculate this ratio, we define a Metropolis function $M(x)= \min\{1, e^{-x}\}$,
and notice that, for a given configuration of the ensemble, the two energy functions have the property
$$ M(U_1-U_0)\exp(-U_0)=M(U_0-U_1)\exp(-U_1).$$
We can then integrate both sides over the configuration space and obtain:
\begin{equation}
    \begin{split}
        \int M(U_1-U_0)\exp(-U_0) dq_1\ldots dq_N & = \int M(U_0-U_1)\exp(-U_1) dq_1\ldots dq_N = \\
        Z_0\frac{\int M(U_1-U_0)\exp(-U_0) dq_1\ldots dq_N}{Z_0}& = Z_1 \frac{\int M(U_0-U_1)\exp(-U_1) dq_1\ldots dq_N}{Z_1}=\\
        Z_0\langle  M(U_1-U_0)\rangle_0 & = Z_1\langle  M(U_0-U_1)\rangle_1 \\
    \end{split}
\end{equation}
So that we obtain the ratio as:
\begin{equation}
 \frac{Z_0}{Z_1}=\frac{\langle  M(U_0-U_1)\rangle_1}{\langle  M(U_1-U_0)\rangle_0}.
\label{eq:ratio}
\end{equation}
The quantities on the right-hand side can be calculated by direct Monte Carlo sampling if $|U_1-U_0|$ is
not too big. One way around if this difference is big is to define intermediate energy functions $U_i$ to
create a chain of highly overlapping ensembles between the two target ensembles. This is especially simple to
do in our particular case of estimating the error rate of QEC codes, as we will see below.

\subsection{Splitting method for estimating QEC thresholds}

An application of the splitting method for the calculation of quantum error correction thresholds can be found in \cite{beverland2025fail}.
In the calculation of QEC thresholds, each energy function corresponds to a physical error rate $p$ for which
we want to find the logical error rate $P(p)$. We begin with a high enough $p_0$ for which we can use regular
Monte Carlo methods to obtain a reliable estimate $P(p_0)$, and obtain $P(p)$ by multiplying over a chain of 
$t$ physical error rates $p_0>p_1>\ldots>p_{t-1}>p$:
$$ P(p) = P(p_0)\cdot\frac{P(p_1)}{P(p_0)}\cdot\ldots\cdot\frac{P(p)}{P(p_{t-1})}$$
If the circuit we are testing has $N$ possible error locations, then a bitstring $E$ of errors has a probability
$\pi_j=p_j^{|E|}(1-p_j)^{N-|E|}$ to occur. Denoting the set of all bitstrings that lead to a logical fault by $\mathcal{F}$,
the probability of a bitstring $E$ given the presence of a logical fault is $\pi_j(E|\mathcal{F})=\pi_j(E)/P(p_j)$.
Then, any function of error configurations can be defined by:
$$\langle f\rangle_j = \sum_{E\in\mathcal{F}} \pi_j(E|\mathcal{F}) f(E)$$
Which we can estimate using $T$ Monte Carlo samples by 
$$\langle \hat{f}\rangle_j = \frac{1}{T}\sum_{\alpha=1}{T}f(E_\alpha)$$
provided that the distribution of $E_\alpha$ approaches $\pi_j(E|\mathcal{F})$.
We can obtain such a set $\{E_\alpha\}_j$ by performing metropolis sampling over error configurations, beginning with
a bitstring $E_0$ that is known to cause a logical fault, and given a last known failing configuration $E_i$,
adding a new bistring $E_{i+1}$ that differs from $E_i$ on
a single error location with probability $q = \min{1, \frac{\pi_j(E_i)}{\pi_j(E_{i+1})}}$.
The $E_i$'s will form a Markov chain that will eventually become a stationary distribution on $\mathcal{F}$ that approximates $\pi_j(E|\mathcal{F})$.

With two Markov chains $\mathcal{M}_{j-1}$ and $\mathcal{M}_j$ we can calculate the ratio $r_j=\frac{P(p_j)}{P(p_{j-1})}$ by:
$$r_j = \frac{P(p_j)}{P(p_{j-1})}
= c \cdot \frac{\langle g(c\pi_{j-1}/\pi_j)\rangle_{j-1}}
{\langle g(c^{-1}\pi_j/\pi_{j-1})\rangle_j}$$
for some constant $c$ whose optimal value is $r_j$, which typically we arrive at by iterating the equation above 3 times.
The algorithm is detailed below.









